\begin{abstract}
    学位论文应有摘要。
    摘要应具有独立性和自含性，即不阅读论文的全文，就能获得必要的信息。
    摘要的内容应包含与论文等同量的主要信息，供读者确定有无必要阅读全文，也可供二次文献采用。
    摘要宜说明研究工作目的、方法、结果和结论等，重点是结果和结论。
    摘要的撰写应符合GB/T 6447的规定。

    以中文撰写的学位论文，中文摘要一般字数为300字~1000字，如遇特殊需要字数可略多。
    摘要中宜尽量避免采用图、表、化学结构式、非公知公用的符号和术语。

    每篇论文应选取3个~8个关键词，用显著的字符另起一行，排在摘要的下方。
    关键词应体现论文特色，选用专业术语及词汇，在论文中有明确的出处。关键词的撰写应符合CY/T 173的规定。

    为便于国际交流，宜有英文摘要，另起一页，置于中文摘要页的后面。
    英文摘要的内容宜与中文摘要对应。
    英文摘要的下方应标注与中文关健词对应的英文关键词。

    以外文撰写的学位论文，应以同种语言撰写外文摘要。
    除外文摘要外，应有中文摘要，置于外文摘要页的后面，并标注中文关键词。
\end{abstract}

\begin{keywords}
    关键词1；关键词2；关键词3；关键词4；关键词5
\end{keywords}

\begin{abstract}[en]
    Abstract content corresponds to Chinese abstract.
    Generally, it should have no less than 300 English content words, and the length should be one page.
    If necessary, the word count can be slightly longer.
\end{abstract}

\begin{keywords}[en]
    Keyword 1; Keyword 2; Keyword 3; Keyword 4; Keyword 5
\end{keywords}
